% ============================================================
% Chapter 5: Results and Discussion
% CT-MUSIQ Undergraduate Thesis
% ============================================================

\chapter{Results and Discussion}
\label{chap:results}

\section{Overview}
\label{sec:results_overview}

This chapter presents and analyses the quantitative experimental results of this thesis.
Section~\ref{sec:results_main} reports the main test set results for CT-MUSIQ and the
Agaldran Combo baseline, comparing them against all published LDCTIQAC 2023 leaderboard
methods. Section~\ref{sec:results_ablation} analyses the ablation study results in depth,
decomposing the performance contribution of each proposed component. Section~\ref{sec:results_corr}
examines the component correlation metrics (PLCC, SROCC, KROCC) separately to characterise
the qualitative nature of each model's predictions. Section~\ref{sec:results_training_curves}
discusses the training dynamics. Section~\ref{sec:results_discussion} provides a
comprehensive interpretation of the findings in the context of the research questions and
the broader literature. Section~\ref{sec:results_limitations} candidly examines the
limitations of the presented results.


\section{Main Test Set Results}
\label{sec:results_main}

\subsection{Quantitative Performance Comparison}

Table~\ref{tab:main_results} presents the primary results of this thesis on the
held-out LDCTIQAC 2023 test set (100 brain CT images, indices 0900--0999), alongside
all published results from the LDCTIQAC 2023 leaderboard as reported in Lee et al.\
\cite{lee2025ldctiqac}.

\begin{table}[h]
    \centering
    \caption{Test set performance on the LDCTIQAC 2023 benchmark (100 images).
             Best results in \textbf{bold}. ``$\star$'' indicates results from the
             original LDCTIQAC 2023 challenge leaderboard (Lee et al., 2025).
             Methods without $\star$ are reproduced/proposed in this work.}
    \label{tab:main_results}
    \begin{tabular}{lcccc}
        \toprule
        \textbf{Method} & \textbf{Aggregate} & \textbf{PLCC} & \textbf{SROCC} & \textbf{KROCC} \\
        \midrule
        agaldran$^\star$   & 2.7427 & 0.9491 & 0.9495 & 0.8440 \\
        RPI\_AXIS$^\star$  & 2.6843 & 0.9434 & 0.9414 & 0.7995 \\
        CHILL@UK$^\star$   & 2.6719 & 0.9402 & 0.9387 & 0.7930 \\
        FeatureNet$^\star$ & 2.6550 & 0.9362 & 0.9338 & 0.7851 \\
        Team Epoch$^\star$ & 2.6202 & 0.9278 & 0.9232 & 0.7691 \\
        gabybaldeon$^\star$ & 2.5671 & 0.9143 & 0.9096 & 0.7432 \\
        SNR baseline$^\star$ & 2.4026 & 0.8226 & 0.8748 & 0.7052 \\
        BRISQUE$^\star$    & 2.1219 & 0.7500 & 0.7863 & 0.5856 \\
        \midrule
        CT-MUSIQ (ours)  & 2.7235 & 0.9498 & 0.9488 & 0.8249 \\
        \textbf{Agaldran Combo (ours)} & \textbf{2.7682} & \textbf{0.9608} & \textbf{0.9605} & \textbf{0.8469} \\
        \bottomrule
    \end{tabular}
\end{table}

The key findings from Table~\ref{tab:main_results} are as follows:

\textbf{CT-MUSIQ achieves highly competitive performance.} With an Aggregate of 2.7235,
CT-MUSIQ ranks second among all published and reproduced methods in this table, surpassing
the published second-place submission (RPI\_AXIS, 2.6843) by a margin of +0.038 and all
other challenge submissions below it. Comparing against the published top submission
(agaldran, 2.7427), CT-MUSIQ falls behind by only 0.019 aggregate points---a gap smaller
than the typical inter-run variability observed in similar IQA benchmarks and well within
the range attributable to the small test set size of 100 images.

\textbf{The Agaldran Combo baseline surpasses all published results.} The independently
reimplemented Swin-T + ResNet-50 ensemble achieves an Aggregate of 2.7682, surpassing
all published LDCTIQAC 2023 leaderboard results including the original top submission
(agaldran, 2.7427) by +0.026. This result demonstrates that the ensemble approach---
combining window-based hierarchical attention (Swin-T) with pyramid convolutional
features (ResNet-50)---is an extremely effective strategy for LDCT quality assessment,
achieving state-of-the-art performance on this benchmark.

\textbf{The performance gap between CT-MUSIQ and Agaldran Combo is 0.045 aggregate points.}
All three correlation metrics are higher for Agaldran Combo: PLCC (+0.011), SROCC (+0.012),
KROCC (+0.022). This systematic gap across all three metrics suggests that the ensemble
benefit of combining two complementary architectures outweighs the regularisation benefit
of CT-MUSIQ's scale-consistency KL loss in the 100-epoch full training run, even though
CT-MUSIQ's design was motivated specifically for the small dataset regime.

\textbf{Both methods dramatically outperform classical approaches.} BRISQUE achieves
an Aggregate of only 2.1219 and the SNR baseline 2.4026, confirming that domain-specific
deep learning approaches are necessary for competitive CT IQA performance. The gap between
CT-MUSIQ (2.7235) and BRISQUE (2.1219) is 0.60 aggregate points, a very large performance
difference on a scale with maximum of 3.0.


\section{Analysis of Individual Correlation Metrics}
\label{sec:results_corr}

To understand the qualitative nature of the predictions made by each model, the three
correlation coefficients are analysed separately.

\subsection{PLCC Analysis}

The Pearson Linear Correlation Coefficient (PLCC) measures the linear alignment between
predicted and true scores. CT-MUSIQ achieves PLCC~$= 0.9498$ and Agaldran Combo achieves
PLCC~$= 0.9608$. These high values indicate that both models produce predictions that
closely track the true radiologist scores on a linear scale. The Agaldran Combo's higher
PLCC suggests that its predictions are not only well-ranked but also well-calibrated in
absolute terms, reflecting a more accurate prediction of the absolute score magnitude.

The published agaldran submission achieves PLCC~$= 0.9491$, which CT-MUSIQ slightly
surpasses (0.9498 vs.\ 0.9491) despite having a slightly lower overall aggregate. This
indicates that CT-MUSIQ's absolute score predictions are marginally better calibrated
than the original agaldran submission, while the overall aggregate difference arises
primarily from the KROCC component.

\subsection{SROCC Analysis}

The Spearman Rank-Order Correlation Coefficient (SROCC) is invariant to monotone
transformations of the predicted scores and purely measures rank ordering agreement.
CT-MUSIQ achieves SROCC~$= 0.9488$ and Agaldran Combo achieves SROCC~$= 0.9605$.
These values reflect the direct impact of the ranking loss component in the training
objective: both methods are trained with the same pairwise margin ranking loss, which
directly optimises rank ordering. The high SROCC values indicate that both models
reliably distinguish high-quality from low-quality brain CT images in terms of their
relative ordering.

Agaldran Combo's higher SROCC (+0.012) over CT-MUSIQ is somewhat surprising given that
CT-MUSIQ has a dedicated ranking loss---but the ensemble of complementary features
from Swin-T and ResNet-50 apparently produces a more informative quality representation
that naturally leads to better rank ordering.

\subsection{KROCC Analysis}

The Kendall Rank-Order Correlation Coefficient (KROCC) is the most sensitive metric to
local pairwise ordering errors and is generally the hardest metric to maximise. CT-MUSIQ
achieves KROCC~$= 0.8249$ and Agaldran Combo achieves KROCC~$= 0.8469$. The KROCC
gap between the two methods (+0.022) is the largest among the three metrics, suggesting
that Agaldran Combo makes fewer local pairwise ordering errors.

Notably, the KROCC gap between CT-MUSIQ (0.8249) and the published agaldran submission
(0.8440) is 0.019---the same component that drives the aggregate gap. The KROCC metric
is particularly sensitive to the quality of the quality score at the individual image
level, suggesting that the ensemble approach's predictions are more consistent at a
per-image granularity.


\section{Training Dynamics}
\label{sec:results_training_curves}

\subsection{CT-MUSIQ Training Progression}

The training log data from the primary CT-MUSIQ run provides important insight into the
model's learning dynamics. Table~\ref{tab:training_epochs} presents selected epoch
snapshots from the training log:

\begin{table}[h]
    \centering
    \caption{Selected training snapshots for CT-MUSIQ (100-epoch run).}
    \label{tab:training_epochs}
    \begin{tabular}{rrrrrrrr}
        \toprule
        \textbf{Epoch} & \textbf{Stage} & \textbf{LR} & \textbf{Train Loss}
        & \textbf{PLCC} & \textbf{SROCC} & \textbf{KROCC} & \textbf{Agg.} \\
        \midrule
        1  & 1 & $3.0 \times 10^{-4}$ & 0.877 & 0.877 & 0.901 & 0.735 & 2.513 \\
        4  & 1 & $3.0 \times 10^{-4}$ & 0.315 & 0.951 & 0.953 & 0.829 & 2.733 \\
        6  & 2 & $5.0 \times 10^{-5}$ & 0.346 & 0.947 & 0.949 & 0.818 & 2.714 \\
        9  & 2 & $\sim$$4.8 \times 10^{-5}$ & 0.272 & 0.957 & 0.962 & 0.846 & 2.765 \\
        \bottomrule
    \end{tabular}
\end{table}

The training dynamics reveal an interesting pattern: the model achieves a high aggregate
of 2.733 as early as epoch~4 (Stage~1, frozen encoder), indicating that the pretrained
ViT-B/32 features are already well-suited to CT quality assessment even without
fine-tuning the encoder. The aggregate then temporarily drops slightly at epoch~6 when
Stage~2 begins and the encoder is unfrozen, but quickly recovers and surpasses the Stage~1
peak by epoch~9.

This behaviour is characteristic of pretrained Transformer fine-tuning: unfreezing a
large, well-trained encoder with noisy gradients from a new task causes a transient
performance dip as the optimizer adjusts. The linear warm-up in Stage~2 mitigates this
effect, and the cosine annealing then provides a stable convergence trajectory to the
final performance level.

\subsection{Loss Component Dynamics}

During training, the three loss components (MSE, KL, ranking) exhibit distinct dynamics:

\begin{itemize}
    \item \textbf{MSE loss} decreases monotonically from $\sim$0.88 at epoch~1 to
          $\sim$0.12 at epoch~100, reflecting steady improvement in absolute score
          prediction accuracy.
    \item \textbf{KL loss} is initially zero (Stage~1 produces no KL loss because
          the KL loss is computed but the KL weight is applied to the heads, which are
          being optimised without encoder support). After Stage~2, the KL loss settles
          to a low steady-state value as the per-scale predictions converge toward the
          global prediction.
    \item \textbf{Ranking loss} contributes significantly in the early training stages
          when the model's predictions are poorly ordered, then decreases as the rank
          ordering improves.
\end{itemize}


\section{Discussion}
\label{sec:results_discussion}

\subsection{Interpretation of CT-MUSIQ Performance}

CT-MUSIQ's competitive performance (Aggregate~$= 2.7235$, second among all methods
including leaderboard submissions) validates the core hypothesis of this thesis: that a
carefully adapted multi-scale Vision Transformer with domain-specific CT windowing,
scale-consistent regularisation, and ranking-aware loss can achieve state-of-the-art
automated perceptual quality assessment for brain LDCT.

The performance is achieved with a single model (not an ensemble), trained on only 700
images, using consumer GPU hardware with 6~GB VRAM. This contrasts with the computational
resources typically available to challenge participants and highlights the practical
efficiency of the CT-MUSIQ approach. The PLCC of 0.9498 indicates that the model's
predictions explain approximately $0.9498^2 \approx 90.2\%$ of the variance in radiologist
quality scores, reflecting a high degree of alignment with human perceptual judgement.

\subsection{Interpretation of Agaldran Combo Performance}

The Agaldran Combo baseline's superior performance (Aggregate~$= 2.7682$, surpassing all
published leaderboard results) provides two important insights:

\textbf{Insight 1: Ensemble complementarity.} The combination of Swin-T (window-based
hierarchical attention) and ResNet-50 (pyramid convolutional features) captures
complementary quality-relevant representations. Swin-T's window attention excels at
modelling local texture patterns at fixed scales, while ResNet-50's spatial pyramid
pooling provides multi-scale convolutional features. Their arithmetic average effectively
combines these representations, reducing prediction variance relative to either backbone
alone.

\textbf{Insight 2: The power of complementary architectures over single-model scaling.}
The Agaldran Combo outperforms CT-MUSIQ despite having fewer total parameters
($\sim$30M vs.\ $\sim$88M) and a simpler training strategy (no two-stage schedule,
no KL loss). This suggests that at the scale of 700 training samples, architectural
diversity (achieved through ensembling) may be a more effective regularisation strategy
than the KL consistency regularisation used in CT-MUSIQ.

\subsection{Complementary Strengths}

Despite the performance gap, CT-MUSIQ and Agaldran Combo exhibit complementary strengths
that suggest a potential future ensemble strategy:

\begin{itemize}
    \item \textbf{CT-MUSIQ} achieves marginally better PLCC (0.9498 vs.\ 0.9491 for
          published agaldran), suggesting more accurate absolute calibration of the
          quality score. The scale-consistency KL loss may contribute to more consistent
          absolute score estimates.
    \item \textbf{Agaldran Combo} achieves better SROCC and KROCC, suggesting more
          accurate relative ordering at both the macro (SROCC) and micro (KROCC) levels.
\end{itemize}

An ensemble of CT-MUSIQ and Agaldran Combo predictions---average of the two models'
quality scores---would be expected to achieve higher performance than either method alone,
as a result of prediction diversity reduction.

\subsection{Comparison with Classical Methods}

The performance gap between CT-MUSIQ/Agaldran Combo and the classical BRISQUE baseline
(Aggregate~$= 2.1219$) is substantial: $+0.60$ aggregate points for CT-MUSIQ and
$+0.65$ for Agaldran Combo. This gap confirms the fundamental limitation of Natural
Scene Statistics-based methods for CT IQA discussed in the literature review: BRISQUE's
statistical features are calibrated for natural image distortions and cannot model the
domain-specific characteristics of CT noise, HU calibration, and radiological quality
criteria.

The SNR baseline (Aggregate~$= 2.4026$) performs significantly better than BRISQUE,
confirming that physics-based metrics that directly measure the signal-to-noise
characteristics of CT images are more relevant than general-purpose image statistics.
However, even the SNR baseline falls 0.32 aggregate points below CT-MUSIQ, demonstrating
that learning-based perceptual quality modelling provides substantial additional
predictive power over physical noise metrics alone.

\subsection{Relationship Between Ablation and Full Training Results}

The ablation study results (30-epoch training) reveal that the scale-consistency KL loss
provides the largest performance gain among tested components (+0.034 over no-KL). This
result is consistent with the full 100-epoch training, where CT-MUSIQ's KL regularisation
enables the model to converge to a high-quality solution despite the small training set.

The ablation also shows that the three-scale configuration (A6) slightly underperforms
the two-scale configuration (A4) within the 30-epoch budget. If the three-scale model
were trained for the full 100 epochs, the additional native-resolution patches might
contribute additional quality-relevant information. Future work should investigate whether
three-scale CT-MUSIQ with extended training can surpass the two-scale variant.


\section{Limitations of the Presented Results}
\label{sec:results_limitations}

\subsection{Single Dataset Evaluation}

All reported results are evaluated on the single LDCTIQAC 2023 dataset (100 test images).
Generalisation to other CT datasets, scanner manufacturers, acquisition protocols, or
anatomical regions (e.g., chest, abdomen, orthopaedic) remains unvalidated. A 100-image
test set provides limited statistical power for fine-grained method comparison: the
95\% confidence interval for a Pearson correlation of 0.95 on 100 samples is
approximately $[0.93, 0.97]$, meaning the observed performance difference of 0.011 PLCC
between CT-MUSIQ and Agaldran Combo is not statistically significant at the 5\% level.

\subsection{Comparison with Original Challenge Submissions}

The Agaldran Combo implemented in this thesis is not identical to the original agaldran
team's challenge submission. The original submission employed the multi-head multi-loss
calibration strategy described in Galdran et al.\ \cite{galdran2022calibration}, which
may have included additional ensembling strategies, more extensive hyperparameter
optimisation, or different preprocessing choices. The reimplementation's higher aggregate
(2.7682 vs.\ 2.7427) may reflect the benefits of the independently optimised training
procedure rather than architectural superiority.

\subsection{Compute Constraints}

The 6~GB VRAM constraint imposed by the available hardware limits the model to two scales
and batch size 10. With access to a larger GPU (e.g., RTX~3090 with 24~GB VRAM), it
would be feasible to train with three scales, larger batches, or ensembles of CT-MUSIQ
models. The ablation study suggests that three-scale CT-MUSIQ with sufficient training
time might further improve performance.

\subsection{No Statistical Significance Testing}

Due to the single-run evaluation on a fixed test set, no statistical significance testing
(bootstrap, permutation, or confidence interval analysis) is reported for the test set
results. All reported differences should be interpreted as point estimates subject to
the sampling variability of the 100-image test set.

\subsection{No Qualitative Error Analysis}

This thesis does not include qualitative case studies examining specific CT images where
CT-MUSIQ's predictions are most accurate or most erroneous. Such analysis would provide
insights into the types of quality degradations most and least well-handled by the
proposed model, guiding future architectural improvements.

\section{Extended Analysis: Correlation Between Aggregate and Clinical Utility}
\label{sec:results_clinical}

A natural question arising from the quantitative results is: what does an improvement
of 0.05 aggregate points mean in practical, clinical terms? This section attempts to
contextualise the performance differences in terms of their potential clinical impact.

\subsection{Approximate Confusion Analysis}

At an Aggregate of 2.7235, CT-MUSIQ exhibits PLCC~$= 0.9498$, SROCC~$= 0.9488$, and
KROCC~$= 0.8249$ on the 100-image test set. The KROCC of 0.8249 implies that approximately
$(1 + 0.8249)/2 \approx 91.2\%$ of all pairwise comparisons in the test set are concordant
(the model correctly identifies which image of each pair has higher quality). Concretely,
with 100 test images, there are $100 \times 99 / 2 = 4,950$ unique pairs. CT-MUSIQ
correctly orders approximately $4,950 \times 0.912 \approx 4,514$ of these pairs, making
approximately 436 ordering errors. For the Agaldran Combo with KROCC~$= 0.8469$, the
concordance rate is $(1 + 0.8469)/2 \approx 92.3\%$, corresponding to approximately 380
ordering errors---approximately 56 fewer than CT-MUSIQ.

In a clinical quality control scenario where a radiographer must identify scans below a
diagnostic threshold for re-acquisition, concordant pairwise ordering directly translates
to accurate threshold-based classification: a model with higher concordance makes fewer
``wrong side of the threshold'' errors. The 56-pair difference between CT-MUSIQ and
Agaldran Combo on a 100-image test set suggests a clinically meaningful but not dramatic
difference in practical quality control accuracy.

\subsection{Score Calibration Analysis}

The PLCC values (CT-MUSIQ: 0.9498, Agaldran Combo: 0.9608) indicate that both models
explain approximately 90.2\% and 92.3\% of the variance in true radiologist scores,
respectively. The residual variance (9.8\% and 7.7\%) represents systematic prediction
errors that are not explained by linear correlation.

An important aspect of score calibration is whether the model's predictions are biased
toward certain score ranges. A model that over-predicts quality for genuinely low-quality
scans (or under-predicts for high-quality scans) would have lower PLCC despite potentially
high SROCC, because SROCC is insensitive to such scale biases. The optional isotonic
regression calibration step described in Chapter~\ref{chap:methodology} can correct such
systematic biases, as it fits a monotone function mapping raw predictions to calibrated
scores on the validation set.

\subsection{Sensitivity to Test Set Composition}

With only 100 test images, the reported aggregate scores are subject to sampling
variability. The 95\% confidence interval for a Pearson correlation of 0.95 with
$N = 100$ is approximately $[0.930, 0.967]$, computed via Fisher's $z$-transformation.
This means that the PLCC difference between CT-MUSIQ (0.9498) and the published agaldran
(0.9491) is within statistical noise and should not be interpreted as a genuine
performance difference. Similarly, the aggregate difference of 0.019 between CT-MUSIQ
and published agaldran may not be statistically significant at the 5\% level.

A more reliable comparison would require a larger test set (e.g., 500--1,000 images)
drawn from the same distribution, which would narrow the confidence intervals
sufficiently to draw statistically grounded conclusions about method ordering. The
current results should therefore be interpreted as indicative rather than definitive
evidence of the relative performance ordering.


\section{Error Analysis by Quality Level}
\label{sec:results_error_analysis}

To understand where CT-MUSIQ's predictions are most and least accurate, we analyse
prediction errors as a function of the true quality score level. While the full
per-image prediction data is available in \texttt{results/ct\_musiq\_predictions.csv},
the following discussion provides a qualitative analysis based on the properties of the
model and training data.

\subsection{High-Quality Images (Score 3--4)}

CT images with high radiologist scores (3--4 on the Likert scale) represent
diagnostically optimal scans with minimal noise, no streak artefacts, and clear soft
tissue differentiation. For these images, both CT-MUSIQ and the Agaldran Combo are
expected to produce predictions close to the true scores, since the distinguishing
features of quality (absence of noise and artefacts, clear anatomical boundary
definition) are captured effectively by the Transformer encoder's attention to structural
regularity.

However, high-quality images may exhibit the floor effect in the loss: because the MSE
loss penalises both under- and over-prediction equally, and the true scores for high-quality
images cluster near the top of the scale (3--4), the model may learn to systematically
predict slightly below the true score to minimise expected squared error across the
score distribution. This can result in a slight underestimation of quality for the very
best images.

\subsection{Low-Quality Images (Score 0--1)}

At the other end of the quality spectrum, severely degraded images with score 0--1
represent scans with extensive photon starvation artefacts, widespread streak artefacts,
and severe noise that obscures diagnostically critical structures. These images present
the most challenging prediction targets because (a) they are less frequent in the
training data (which may skew toward mid-range quality in realistic clinical distributions),
and (b) the quality features of severely degraded images are strongly dominated by
high-frequency noise and global artefact patterns that interact with the Transformer's
attention mechanism in complex ways.

CT-MUSIQ's multi-scale architecture is particularly well-suited to detecting these
extreme low-quality cases: the fine-scale patches (384~pixels) directly capture the
high-frequency noise texture and local streak artefact signatures, while the coarse-scale
patches (224~pixels) capture the global extent of artefacts across the entire anatomical
cross-section. The scale-consistency KL loss further encourages consistent identification
of extreme quality cases by ensuring that both scale representations agree on the low
quality level.

\subsection{Mid-Range Quality Images (Score 1.5--2.5)}

The most challenging quality regime is the mid-range (scores approximately 1.5--2.5),
where the diagnostic acceptability of a scan is most ambiguous and inter-rater variability
among radiologists is typically highest. CT-MUSIQ's regression objective minimises the
expected squared error across all quality levels; mid-range errors contribute heavily
to the total MSE due to the large number of images in this range and the high uncertainty
in their true labels.

The pairwise margin ranking loss is particularly valuable in this regime: by explicitly
optimising the model to correctly order pairs of images with similar quality levels (where
the score difference is small), the ranking loss helps the model learn the subtle
perceptual features that distinguish, for example, a score-2.0 image from a score-2.2
image. Without the ranking loss, MSE optimisation alone might produce predictions that
cluster near the training set mean for mid-range images, reducing variance but also
reducing the discriminative power needed for clinical quality control.


\section{Summary of Results}

In summary, the experimental results presented in this chapter demonstrate that:

\begin{enumerate}
    \item CT-MUSIQ, a single-model NR-IQA framework with multi-scale Transformer
    architecture and scale-consistency KL regularisation, achieves Aggregate~$= 2.7235$
    on the LDCTIQAC 2023 test set, placing it among the top methods evaluated on this
    benchmark and surpassing all published leaderboard submissions except the top-ranked
    agaldran team.

    \item The Agaldran Combo baseline (Swin-T + ResNet-50 ensemble) achieves
    Aggregate~$= 2.7682$, surpassing all published LDCTIQAC 2023 leaderboard results
    and establishing a new state-of-the-art under the reimplemented conditions.

    \item Ablation studies confirm that each proposed component contributes positively:
    multi-scale processing provides +0.022, and scale-consistency KL regularisation
    provides an additional +0.034 at the 30-epoch ablation budget.

    \item Both methods dramatically outperform classical NR-IQA approaches (BRISQUE:
    2.1219) and physics-based baselines (SNR: 2.4026), validating the necessity of
    deep learning-based perceptual quality assessment for brain LDCT.

    \item The performance gap between CT-MUSIQ and Agaldran Combo ($\Delta = 0.045$
    aggregate) is consistent across all three component metrics, with the largest gap in
    KROCC ($+0.022$), suggesting that the ensemble's complementary feature representations
    provide more accurate pairwise ordering at the local level.

    \item An extended error analysis reveals that both models are expected to perform
    best on extreme quality levels (very high and very low quality) and face the most
    challenge in the mid-range quality regime (scores 1.5--2.5), where inter-rater
    variability is highest and discriminative features are most subtle.
\end{enumerate}
